% =====================================================================
% Sekcja 6 - Dyskusja
% =====================================================================

\section{Dyskusja}
\label{sec:dyskusja}

Zebrane wyniki układają się w kilka spójnych wątków, które warto
omówić razem - bo dopiero zestawione ze sobą i z literaturą z
sekcji~\ref{sec:tlo} mówią coś więcej niż pojedyncze współczynniki
korelacji.

Pierwszy wątek dotyczy wyboru metryki. We wszystkich ustawieniach -
od pilotażowego modelu 0,5B po agenta 1,5B - średnia entropia
tokenowa i top-1 margin okazują się lepszymi predyktorami
poprawności niż samo perplexity. Różnica jest niewielka, ale
konsekwentna, i ma dość naturalne wytłumaczenie: entropia i margin
korzystają z całego rozkładu prawdopodobieństw, a perplexity tylko
z prawdopodobieństw tokenów, które model faktycznie wybrał. Model
może wybrać token z wysokim prawdopodobieństwem, mając tuż za nim
silną alternatywę - perplexity tego nie zauważy, a margin i entropia
tak. Sugeruje to, że w dalszych pracach nad pewnością agentów warto
sięgać po metryki rozkładowe, a nie ograniczać się do klasycznego PPL.

Drugi wątek to lokalizacja pomiaru wewnątrz agenta. Sygnał nie jest
rozłożony równomiernie po trajektorii: najsilniej koreluje z
poprawnością pewność na segmentach \texttt{Action}, słabiej na całej
trajektorii, a na samym segmencie \texttt{Answer} znika niemal
zupełnie. Wynika to z mechaniki ReAct - w momencie formułowania
finalnej odpowiedzi wynik liczbowy jest już w kontekście, więc model
go po prostu przepisuje i jest pewny niezależnie od tego, czy droga
do tego wyniku była poprawna. Pewność niesie informację tam, gdzie
model rzeczywiście podejmuje decyzję (co policzyć), a nie tam, gdzie
tylko odtwarza gotowy rezultat. To ustalenie jest zbieżne z
obserwacją Fanga i in.~\cite{fang2025longppl}, że uśrednianie po
całej długiej sekwencji rozcieńcza sygnał - tyle że tutaj rolę
,,nieistotnych'' tokenów grają nie spójniki, lecz cała końcowa
faza kopiowania odpowiedzi.

Trzeci, najważniejszy wątek dotyczy różnicy między korelacją a
kalibracją (sekcja~\ref{sec:kalibracja}). Perplexity porządkuje
pytania sensownie, ale jako bezwzględna miara pewności jest źle
wykalibrowane - model jest systematycznie zbyt pewny siebie, i to
tym bardziej, im trudniejsze ustawienie. Rozstrzyga to pozornie
sprzeczne stanowiska z literatury. Gonen i
in.~\cite{gonen2023demystifying} mają rację, że niższe PPL wiąże się
z lepszymi wynikami - kierunek jest poprawny. Zastrzeżenia Wanga i
in.~\cite{wang2022perplexity} co do wiarygodności PPL też są jednak
zasadne, jeśli traktować je jako uwagę o skali, a nie o porządku.
Te dwa stanowiska nie są sprzeczne - opisują dwie różne własności
tej samej liczby.

Płynie stąd praktyczny wniosek dla prac, które wbudowują pewność w
pętlę decyzyjną agenta (sekcja~\ref{sec:agent-uncertainty}). ReDAct
\cite{redact2026} i pokrewne podejścia opierają się na progowaniu
surowej pewności, a wyniki tej pracy pokazują, że taki próg nie
przenosi się między zadaniami bez rekalibracji: ta sama wartość
liczbowa oznacza inną faktyczną szansę poprawności w czystym CoT i
w agencie. Sam pomysł sterowania niepewnością pozostaje sensowny,
ale wymaga albo kalibracji na danych docelowych, albo oparcia decyzji
na względnym uporządkowaniu, a nie na sztywnym progu.

Na koniec warto odnotować, że dwa główne wzorce - przewaga metryk
rozkładowych oraz przeniesienie sygnału na segmenty \texttt{Action} -
utrzymują się przy przejściu z modelu 0,5B na 1,5B oraz między
ziarnami losowości. Przy skali tego projektu nie jest to dowód
ogólności, ale sugeruje, że obserwacje nie są wyłącznie artefaktem
jednego modelu czy jednego losowania danych.
